\section{Two marked bonds on tilted cylinders}\label{c7:marked}

The length-square estimate in \eqref{j:polygon-square} will follow by
summing the weight of polygons through pairs of specified bonds. We
first represent that positive sum by a finite matrix element. The
quantitative estimate allows the two intervals between the bonds to
have different lengths.

\subsection{The physical sum and the estimate to prove}

Put $\tau=e^{i\pi/3}$. After a lattice rotation or reflection, a
\emph{permitted cut} is a periodic path of tessellation sides whose
steps are $1$ and $\tau$. Translating it by the time step $\tau^2$
gives rows of rhombi: the spatial projection transverse to $\tau^2$
is strictly increasing, so these rows tile the plane. A period has
$M$ sides and a nonzero lattice translation $P$. We call a site
horizontal when its cut step is $1$, and mark the honeycomb bond
crossing each of two horizontal sites $i,j$. Choose the seam just
before $i$. The period then consists of consecutive nonempty lists
$A,B$, with $i$ first in $A$ and $j$ first in $B$; write
$A'=A\setminus\{i\}$ and $B'=B\setminus\{j\}$.

On the cylinder modulo $P$, let $H^{\rm poly}$ be the increasing
closed-rim height limit of
\[
 \sum_{\substack{\lambda\ {\rm a\ single\ simple\ unoriented\ cycle}\\
                   \lambda\ {\rm uses\ both\ marked\ bonds}}}
 x^{|\lambda|}.
\]
The sum includes contractible and noncontractible cycles. The
proposition below proves that the limit is finite.

The associated physical site variables have shifted-to-horizontal
ratio $q=e^{i\pi/4}$. Indeed the angles of the two rhombi relative
to the time step are $\pi/3$ and $2\pi/3$; at one common physical
row their spectral differences are $\alpha$ and $2\alpha$ for the
shifted and horizontal sites, respectively. At both differences
the scalar weights are exactly the honeycomb weights after splitting
each rhombus into two equilateral triangles. At $2\alpha$ the
weights $u,d$ and $a,b$ are interchanged relative to $\alpha$.

For a generic nonzero site list $X$ on the same two lists, define a
separate algebraic quantity
\begin{equation}\label{j:marked-flip-definition}
 H^{\rm flip}(X)=
 L_- f_i\, i^{\sum_{a\in A'}s_a^{\rm phys}} e_j R_+,\qquad
 f_i=|-\rangle_i\langle+|,\quad e_j=|+\rangle_j\langle-|.
\end{equation}
Here $s_a^{\rm phys}$ is the diagonal site-spin operator, and
$R_+=R_i$, $L_-=L_{-i}$ are the selected vacuum vectors of
Lemma~\ref{j:selected-vacuum-vectors}. Their generic normalization
is the one fixed in Section~\ref{c7:finite}.

\begin{proposition}[Physical two-bond identity]\label{j:marked-physical}
Let the permitted cut have even $M$, and let both marked sites be
horizontal. The positive sum $H^{\rm poly}$ is finite. The generic
matrix element $H^{\rm flip}$ has a regular limit at the associated
physical site list, and at that list
\[
 H^{\rm flip}=H^{\rm poly}.
\]
This holds for every order and every number of the two site types in
the two nonempty lists; the unmarked lists may be empty.
\end{proposition}

\begin{theorem}[Unequal marked intervals]\label{j:unequal}
Use the physical tilted cylinder and horizontal-type marked bonds
just defined. Its two cut intervals have $m+1$ and $n+1$
sites, with $n\ge m$ and $M=m+n+2$ even. In each unmarked list the logarithmic site
coordinates have imaginary parts $-f$ and $1-f$, where $0\le f\le0.6$. Their respective
counts differ from $m(1-f),mf$ in the first unmarked list and from
$n(1-f),nf$ in the second by at most a fixed constant. Let $H^{\rm poly}$
be the sum of $x^{|\lambda|}$ over single unoriented polygons through
the two specified bonds. Uniformly over these lists and their orders, for
every $\epsilon>0$ there are $C_\epsilon$ and $m_\epsilon$ such that
\[
 H^{\rm poly}\le C_\epsilon n^\epsilon m^{-4/3}(n/m)^{1/3}
 \qquad(n\ge m\ge m_\epsilon).
\]
Both $C_\epsilon$ and $m_\epsilon$ may depend on the fixed bound for
the count discrepancies.
\end{theorem}

The physical identity has the broader domain: its proof uses the
selected vacuum vectors, a finite nonvanishing result for their scalar
normalizer, summability on a fixed cylinder, and a two-arc phase
calculation. The count restrictions in Theorem~\ref{j:unequal} enter
later, when Sections~\ref{c7:markedgas}--\ref{c7:mixed} estimate the
horizontal representation. We keep the algebraic and positive sums
separate until Proposition~\ref{j:marked-physical} is proved.

\subsection{The generic flip formula}

Put $U=x_i$, $V=x_j$, and retain arbitrary nonzero values of these
variables. For the unmarked lists $(A',B')$, let $C_s$ denote the
summand in \eqref{j:finite-spin-sum} for $N^2\mathscr G$ at $h=1$,
including its phases. Thus the labels satisfy
$Q=\sum_{A'}s_a=\sum_{B'}s_a$. For a pair $(u,v)$ write
$e_k(u,v)=v-q^k u$.

\begin{proposition}[Generic long-flip formula]\label{j:marked-flip-formula}
For generic nonzero $X$ on consecutive nonempty lists $A,B$,
\begin{equation}\label{c7:eq-L}
\begin{aligned}
 N(X)^2 H^{\rm flip}(X)
 &=cUV\sum_s q^{-6Q}C_s
       \prod_{a\in A'\cup B'}F_{s_a}(U,x_a)F_{s_a}(V,x_a),\\
 (F_+,F_0,F_-)&=(e_5e_6,e_3e_5,e_2e_3),\qquad c=2-\sqrt2.
\end{aligned}
\end{equation}
The two endpoint variables $U,V$ are independent in this identity.
\end{proposition}

\begin{proof}
We first separate the individual spin terms of the unflipped pairing.
For each of two lists, use its ordered partial monodromy with both
auxiliary ends fixed at $+1$ for the first list and at $-1$ for the
second. Within a list these operators commute as the probe varies:
the exchange proof applies because the equal extremal auxiliary pair
has a scalar block. Label basis spins by $s=s^{\rm phys}$ in the
first list and $s=-s^{\rm phys}$ in the second. At fixed total spin,
order configurations by their prefix sums. Each partial monodromy is
triangular, with row $\le$ column, because the intermediate auxiliary
spin stays within its bounds.

The diagonal characters are products of the one-site entries
$r_s(w/x)$, where
$(r_+,r_0,r_-)=(a_{\rm tile},v_{\rm tile},b_{\rm tile})$.
They are distinct jointly at generic variables. Their denominator
roots are $w/x=q^5,q^6$, and their numerator root indices are
\[
D_-=\{0,1\},\quad D_0=\{0,3\},\quad D_+=\{2,3\}.
\]
Apply the joint eigenprojection to its leading coordinate to obtain
$|s\rangle\rangle$ and $\langle\langle s|$, taking the product over
the two lists. Their leading coefficients are one, their ket and bra
supports are respectively $\le s$ and $\ge s$, and they are dual.
Indeed each generic projection is a polynomial in triangular matrices
with exactly one diagonal entry equal to one.

Consider the two amplitudes $N\langle\langle s|R_+$ and $N L_- |s\rangle\rangle$, in total physical charge zero. We list forced factors, omitting phases. In both there is a factor $\prod_a x_a^{s_a^2/2}$. For a same-group pair in order $x,y$, labels $a,b$, the right amplitude (involving $R_+$) has zeros
\begin{itemize}
\item $e_5(x,y)$ unless $(a,b)=(+,-)$;
\item $e_6(x,y)$ unless $(a,b)\in\{(+,0),(0,-),(+,-)\}$.

\end{itemize}
Allow denominator $e_0$ if $a<b$, $e_1$ for $(-,+)$, $e_7$ for $(0,0)$. For the left amplitude use the same right rule with the order reversed. Thus the product of the two rules gives the rational monic part of $p_{ab}$ above up to constants. For a cross pair (either list first), in \textbf{each} amplitude use a factor $e_5$ if $a<b$, $e_3$ if $a>b$, additionally $e_6$ for $(-,+)$, $e_2$ for $(+,-)$. There are no cross denominator factors.

\paragraph{Forced factors and partial spectra.}
The vacuum input here is precisely the asymmetric statement of
Lemma~\ref{j:selected-vacuum-vectors}. For a physical zero-charge
configuration $\sigma$ and $t_a^2=x_a$, it says that
\[
 \frac{N(R_+)_\sigma}{\prod_a t_a^{\sigma_a^2}},
 \qquad
 \frac{N(L_-)_\sigma}{\prod_a t_a^{\sigma_a^2}}
\]
are polynomials in the ordinary site variables, independent of the
square-root signs, of degree at most $|X|-1-\sigma_a^2$ in $x_a$.
For a vacant coordinate, their value at $x_a=0$ is
$\prod_{b\ne a}x_b$ times the corresponding reduced polynomial,
and their coefficient of $x_a^{|X|-1}$ is that reduced polynomial.
With one varied site and generic fixed spectators, the proof also
gives, for the normalized coordinates, an occupied coordinate of
order $O(t_a)$ at zero and $O(t_a^{-1})$ at infinity, and the
reduced vacant coordinate at both limits. These conclusions apply
to $R_+$ and $L_-$ only.

The additional poles of our amplitudes come from the partial
projections. At a generic point of a same-list divisor
$y=q^\delta x$, their spectral coincidences require equality of
$D_a\cup(\delta+D_b)$ modulo 8 at fixed $a+b$. The displayed root
lists give only the following doublets: unequal swapped labels at
$\delta=0$; $(-,+)$ and $(0,0)$ at $\delta=1$; and their reflections
at $\delta=-1$. The spectator labels agree, and there are no
coincidences at other generic divisors. A regular generic linear
combination of probes separates either doublet to first order:
moving the second ratio moves a different multiset of numerator
roots, so the logarithmic derivatives differ. It also separates
the other characters. Thus these poles have order at most one.
Solving in triangular order on supports shows that the ket of the
lower label and bra of the higher have no pole, since those supports
exclude the colliding index. This is exactly the denominator rule
above. Several characters can coincide at an intersection of these
divisors; we make no doublet assertion there.

The singleton limits of the partial basis require a separate
observation. With one site on its own scale, fixed regular probes
have diagonal entry $q^{\pm3s}$ times the generic spectator
character, so a generic probe has a holomorphic simple projection
at that limit. It preserves the varied site's vacancy parity there.
Let $P_a$ act by $(-1)^{\sigma_a^2}$ on the physical site basis.
The leading-coordinate normalization gives
\[
 |s\rangle\rangle(-t_a)
   =(-1)^{s_a^2}P_a|s\rangle\rangle(t_a),\qquad
 \langle\langle s|(-t_a)
   =(-1)^{s_a^2}\langle\langle s|(t_a)P_a.
\]
Together with the coordinate limits of the selected vectors, this
shows that either amplitude divided by $N$ is bounded at both
singleton ends, and is respectively $O(t_a)$ and $O(t_a^{-1})$
there if $s_a=\pm1$. It also shows that extracting
$\prod_a t_a^{s_a^2}$ makes each amplitude rational in the ordinary
variables.

For same-group fusion put a higher site $q^\delta u$ before a lower $u$, $\delta=2,3$. The site intertwiner used above places $R_+$ in its image and intertwines even the partial monodromy there (for adjacent sites). Simplicities of the full and reduced fixed characters hold generically there by the same paired-scale tests, and the partial spectra are still simple by the root lists. Only labels realizing a reduced eigenvalue can occur in $\langle\langle s|R_+$. Comparing numerator lists after denominator cancellation, in (high, low) order these are $(+,-)$ for $\delta=3$, and $(+,0),(0,-),(+,-)$ for $\delta=2$. For instance at $\delta=2$ the two removed poles in numerator indices relative to $u$ are $0,5$. This proves the same-group right zeros. Swaps moving the pair together while preserving their relative order are allowed by regular exchange similarities at generic spectators, mapping the partial label projections along with the fixed lines. The normalizations cause only nonzero finite scalars (all these lines simple with their stated leading coefficients). The identical reasoning in reversed order proves left zeros: at real phases take Hermitian adjoints of monodromies, which reverses order (and conjugates the twist), then extend identities algebraically.

For cross fusion with $x=q^\delta y$, put the high site last of its list and the low first of its list. This uses same-list swaps, and, if the lists have to be interchanged cyclically, a seam change acting by a diagonal twist on the conserved list spin (hence preserving partial eigenspaces), directly by charge conservation. Suppose first the high site is in the first ($+$ convention) list. In the right-amplitude bra its physical spin is $\le a$ by prefix support; that of the low site is $\le-b$. But the vacuum ket in the fusion image uses total spin on the pair equal to 0 at $\delta=3$, or in $[-1,1]$ at $\delta=2$. This excludes $a<b$ or $a<b-1$ respectively. When the high site belongs to the other list negate physical spins in this argument. Apply the same adjacency arrangement in reversed order to the adjoint for the left amplitude. This proves the claimed cross factors, with both signs of shifts by interchanging arguments.

After dividing by the stated rational factors and site half powers,
the remaining rational function has no pole. Indeed every possible
irreducible divisor at finite nonzero values has been checked at a
generic point, and the singleton parity estimate removes the
possible pole at zero. Unique factorization then removes the same
denominators globally, including at intersections of shifted
divisors. The remainder is a polynomial. For a pair, the net degrees
used by the rational factors are $1+st$ in the same list and
$1-st-\mathbf1_{s=t=0}$ across lists. Charge equality, the
singleton bounds, and $\deg_{x_a}N\le |X|-1$ show that this
polynomial is independent of every nonzero-label variable. Its
degree in a zero-label variable is at most the number of zeros in
the opposite list.

It follows that individually for the unflipped two-list problem
\[
N^2\, i^Q(L_- |s\rangle\rangle)(\langle\langle s|R_+) = C_s
\]
(here $C_s$ and $Q$ temporarily refer to those arbitrary lists). Indeed induct on size. If a list has a distinct-label pair, choose a specialization $(x,y)=(u,q^2u)$ on these labels in \emph{ascending label} order (regardless of site order). This imposes a value on a nonzero-label variable, still free in the factored-out portion, without any zero or pole from its forced factors. By full similarity covariance of the left-right projected product (normalization by vacuum overlap 1), we may arrange the high site immediately before the other (same-list reordering here is an identity first at generic coordinates, continued to the regular values on this specialization even if an exchange itself becomes singular). Site fusion now reduces the cleared expression exactly to the reduced individual label expression times the $N^2$ fusion multiplier: the right fixed vector lies in the image, the partial probes intertwine there with simple spectra, the left fixed vector restricts to the reduced one up to normalization (product of normalizing scalars 1 by overlap), and the list spin is unchanged. The product check for surviving summands in the earlier spin formula has precisely this individual reduction. The displayed candidate shares the rational/degree rules (with the extra cross $00$ numerator), so the single evaluation in the nonzero-label variable proves equality. This handles all nonuniform cases. The remaining cases then follow by summing projections and comparing coefficients of $h^Q$ in the known unflipped spin sum: the all-zero term is the only unaccounted one there at $Q=0$, and a uniform charged configuration, if any, uniquely realizes its extremal $Q$ on two equal nonempty lists. This includes the induction base.

\paragraph{Insertion of the two long flips.}

Insert now the long flips between expansions in the partial bases on the \textbf{full} lists. By first-site triangular support the matrix element from ket label $v$ to bra label $u$ is zero unless $v_i=v_j=+$ and $u_i=u_j=-$. In these branches (physical first spins necessarily extremal due to the flips), the sections of basis vectors at their indicated first labels are exactly the corresponding normalized partial basis vectors on the remaining lists: the fixed-first-label triangular diagonal block of the probes is just the reduced probe times that first diagonal factor. Thus remaining labels have to match, say $s$ there, and the matrix element equals $i^Q$. We must evaluate $N^2(L_-|u\rangle\rangle)(\langle\langle v|R_+)$ with these labels. As functions of $U,V$ the forced factors and degree rule determine this up to a spectator factor: multiplying the two tables gives precisely
\[
UV\prod_a F_{s_a}(U,x_a) F_{s_a}(V,x_a)
\]
for their entire dependence (product over unmarked sites). Here the endpoints were first of their own lists.

Put the endpoints now together on one large scale above the spectators, in generic ratio. In these particular amplitudes $|u\rangle\rangle$ and $\langle\langle v|$ have actual first spins fixed by extremality, and reduced-list sections as just noted. Their physical endpoint total charge is zero, so only the diagonal grade block of the vacuum projections is needed. The two-scale argument above isolates the fixed character as a product of endpoint pair and spectator projections, up to the common diagonal similarity: charged endpoint totals are excluded exactly by the regular pair tests used for matched pair scales ($p_j^{S_j}$ for odd total, and the generic scalar $p_j^{\pm1}a_{\rm tile}$ for absolute total two). Neutral blocks see original twists. On endpoint physical spins $(t,-t)$ and given spectator list totals $(Q,-Q)$ the similarity displayed there acts by $q^{3tQ}$. Hence the amplitudes divided by $N$ each pick up $q^{3Q}$ relative to the reduced-list ones in this limit, as well as the endpoint pair components, whose combined multiplier by the explicit vacuum pair formula ($p=i$ on right, $-i$ on left, respectively entries $+,-$) is $cUV/(U+V)^2$. Also
\[
N(X)\sim (UV)^{|X|-2}(U+V) N(X\setminus\{U,V\});
\]
cross rows in the Pfaffian coincide toward the rest. Finally the leading phase in the displayed $F$ product is $q^{6\sum_a s_a}=q^{12Q}$. This computes the spectator factor by the individual unflipped formula and proves \eqref{c7:eq-L}.
\end{proof}

\subsection{Finite nonvanishing at the physical sites}

The physical vacuum argument needs only finite nonvanishing. The
growth of the same scalar with the number of sites will enter later
in Section~\ref{c7:markedgas}.

\begin{lemma}[Orbital ratio and nonvanishing]\label{j:shifted-nonvanishing}
Let $M=2L\ge2$, $w_j=i h_j$ with $h_j\in\mathbb R$ of spread at
most one, and $x_j=e^{\eta w_j}$, where $\eta=\pi/4$. Define
\[
 \mathcal S(X)=\mathcal P(X)\prod_{j<l}C(w_j-w_l),
 \qquad C(r)=\tanh(\pi r/6),
\]
using the confluent value for equal sites. Then $\mathcal S(X)\ne0$
and
\begin{equation}\label{j:shifted-orbital-ratio}
 \frac{N(X)}{\mathcal S(X)}
 =\left(\frac23\right)^L
   \mathbb E_h\exp\!\left(\sum_{a=1}^L\phi(k_a)\right)>0,
 \qquad
 \phi(k)=\log\left(1+\frac{e^{-2k}}{1+e^{-4k}}\right).
\end{equation}
Here $\mathbb E_h$ has ordered density on $0<k_1<\cdots<k_L$
proportional to
\[
 \Delta(k^2)^2\prod_{a=1}^L\frac{k_a}{\sinh3k_a}\ I_h(k),
 \qquad
 I_h(k)=\int_{U(2L)}
     e^{\operatorname{tr}(D_k\mathcal U D_h\mathcal U^*)}\,d\mathcal U,
\]
where $D_k=\operatorname{diag}(k_1,\ldots,k_L,-k_1,\ldots,-k_L)$
and $D_h=\operatorname{diag}(h_1-h_*,\ldots,h_M-h_*)$ for
$h_*=(\max h_j+\min h_j)/2$. The measure on $U(2L)$ is normalized
Haar measure. In particular $N(X)\ne0$. The nonvanishing also holds
after multiplying every $x_j$ by the same nonzero scalar.
\end{lemma}

\begin{proof}
First, each pair factor
$d(x_j,x_l)C(w_j-w_l)/(x_j-x_l)$ is nonzero on this domain,
including its finite nonzero limit at equality. Indeed the roots of
$d(u,v)=u^2+\sqrt2uv+v^2$ have relative arguments
$\pm3\pi/4$, outside the present range $[-\pi/4,\pi/4]$.
Within $|\Im(w_j-w_l)|\le1$, the only zero of $C$ is the simple
zero at equality and $C$ has no pole. This proves the assertion
about $\mathcal S$.

For distinct $h_j$, use the transforms of $B_0'$ and $C'$ in the
homogeneous normalization calculation. For either odd kernel
$B=B_0,C$,
\[
 B(i(h_j-h_l))
 =\int_0^\infty
   \frac{i\sinh((h_j-h_l)k)}{\pi k}\widehat{B'}(k)\,dk.
\]
The $L$th exterior power of the integral of these decomposable
two-forms expresses its Pfaffian as an integral of
$\det[e^{\pm h_j k_a}]$ on ordered positive $k_a$, with the product
of the corresponding weights $\widehat{B'}(k_a)/k_a$ and common
constants. Apply the unitary orbital determinant identity derived
before \eqref{j:boundary-orbital-law}. The Vandermonde of the paired
$k$ list supplies a constant times
$\prod_a k_a\,\Delta(k^2)^2$. Schur's identity gives
$\operatorname{Pf}C(w_j-w_l)=\prod_{j<l}C(w_j-w_l)$, while
\[
 \frac{\widehat{B_0'}(k)}{\widehat{C'}(k)}
 =\frac23 e^{\phi(k)},\qquad
 \widehat{C'}(k)=\frac{6k}{\sinh3k}.
\]
Taking the ratio proves \eqref{j:shifted-orbital-ratio}.

All integrations are absolute on the stated domain:
$\|D_h\|\le1/2$ gives $I_h(k)\le e^{\sum_a k_a}$, whereas
$k/\sinh3k$ decays like a polynomial times $e^{-3k}$.
This bound also permits coincident $h_j$ by continuity in the
orbital formula. The expectation is strictly positive at every
size, proving $N\ne0$. Homogeneity of $N$ and $\mathcal S$ gives
the last assertion.
\end{proof}

\subsection{Summability on a fixed cylinder}

\begin{lemma}[Period exclusion]\label{j:period-excluded}
Fix a nonzero lattice translation $P$ and a honeycomb vertex $v$.
The total critical weight of plane SAWs starting at $v$ whose
visited vertices are distinct modulo $P$ is finite. The empty walk
may be included. The bound may depend on $P$.
\end{lemma}

\begin{proof}
Rotate the lattice and, if necessary, replace $P$ by $-P$ so that
$P=a\tau+b(\tau-1)$ with integers $a,b\ge0$. These choices are
possible because the sixty-degree cones spanned by successive
lattice directions cover the plane. Use horizontal tessellation
lines for these coordinates. Let $a_h$ be the total weight of
first-exit slab crossings from a fixed bottom port to any top port
$h$ layers above, and let $b_h$ be the weight of those crossings
with no intermediate level crossed exactly once. The convex balance
in Section~\ref{j:local}, exhausted horizontally, gives
$a_h\le\csc(\pi/8)$ for every $h$.

Cut a crossing at every intermediate level that it crosses exactly
once. The resulting blocks are irreducible crossings counted by the
$b_h$. Conversely, any word of such blocks concatenates freely:
the first-exit property places successive blocks in disjoint open
slabs, sharing only their joining port. Their weights multiply.
Thus, for $0<t<1$,
\[
 A(t):=\sum_{h\ge1}a_ht^h
   =\sum_{r\ge1}B(t)^r<\infty,\qquad
 B(t):=\sum_{h\ge1}b_ht^h.
\]
It follows that $B(t)<1$ and, by monotone convergence,
$\mu:=\sum_{h\ge1}b_h\le1$.

Among the height-one blocks, choose the word that takes the $a+b$
upward moves whose sum is $P$ and then repeats its first block.
Its length is $r_P=a+b+1$ and its weight is
$p_P=x^{2r_P}>0$: each move traverses two triangles. Its first and
last repeated blocks visit vertices equal modulo $P$. A
period-excluding word must therefore avoid this specified word.
Divide any longer word into disjoint chunks of length $r_P$ and a
last shorter chunk. The sum of weights of words that avoid the
specified word is at most
\[
 \sum_{j=0}^{r_P-1}\mu^j
       \sum_{r\ge0}(\mu^{r_P}-p_P)^r
 \le \frac{r_P}{p_P}<\infty.
\]
The same remains true if a first and a last irreducible block are
left unrestricted, since their total masses are at most $\mu$.

We transfer this bound from ports to vertices. Measure a vertex
span $l$ in units of one third of the layer spacing. Extend a
segment whose endpoints are opposite height extrema to the adjacent
port lines, as in Lemma~\ref{j:extremal-segment}. The extension
adds at most one triangle at either end, at bounded multiplicity.
If the resulting slab has height $h$, then
$3h-l\in\{2,3,4\}$. The sequence of middle irreducible blocks
still belongs to the original period-excluding segment, so it
avoids the specified word. The preceding sum therefore gives
bounds $C_l$ for the weight of such segments from a fixed vertex to
all opposite extremal endpoints, with
\[
 \sum_{l\ge1}C_l<\infty.
\]
We take the maximum over both directions and the two vertex types;
the port extension costs only a fixed factor.

A segment starting at its highest vertex and staying no higher
splits at successive last opposite extrema as described after
Lemma~\ref{j:extremal-segment}. Its positive spans decrease
strictly, apart from a possible equality of the first two. Put
$d_l=C_l/x$, accounting for the shared vertex at a join. Summing
over these span lists costs at most a fixed factor times
\[
 \left(1+\sum_{l\ge1}d_l^2\right)
 \prod_{l\ge1}(1+d_l)<\infty.
\]
The empty segment is included. Finally split a full walk at a
specified highest vertex, for example the first such vertex. To
count the reversed first half without summing over free highest
positions, translate its highest endpoint to one of two fixed
representatives, according to vertex type. Its translated end and
the original fixed start determine the translation. Both halves
are therefore bounded by the preceding endpoint-summed bound.
\end{proof}

\subsection{The phase of two marked arcs}

We calculate phases on the reference flat cylinder with horizontal
cut and angle $\pi/3$, keeping the scalar weights of the physical
tilted rows. The local pictures and phases after the telescoping
removal in Section~\ref{j:local} are unchanged by this choice.
Place the seam at zero and write the endpoint abscissae as
$0<i<j<M$.

Consider a smooth regular embedded lift of one directed arc
$X:[0,1]\to\mathbb R^2$ from $(i,0)$ to $(j+kM,0)$,
$k\in\mathbb Z$, disjoint from all of its nontrivial horizontal
$M$-translates. Its endpoint tangents are vertical with signs
$\zeta_i,\zeta_j\in\{-1,1\}$, upward positive. Use a small
perturbation to make its interior cut and seam-ray crossings
transverse, retaining the absence of interior cut crossings at
seam or endpoint abscissae and its prescribed tangent turns.
The actual reference lattice arcs have this absence before
smoothing. Let $x_z$ and $\epsilon_z$ be the abscissae and signed
directions of the interior cut crossings, and define
\[
 \mu=\sum_{\text{interior crossings }z}\epsilon_z\delta_{x_z}
       +\frac{\zeta_i}{2}\delta_i
       +\frac{\zeta_j}{2}\delta_{j+kM},\qquad
 C(a)=\mu((-\infty,a))+\frac12\mu(\{a\}).
\]
The total mass is zero, since the interior signed crossing total is
$-(\zeta_i+\zeta_j)/2$. In particular $C$ has compact support.
Define the finite sums
\[
 C_0=\sum_{\ell\in\mathbb Z}C(\ell M),\qquad
 C_{\rm per}(a)=\sum_{\ell\in\mathbb Z}C(a+\ell M)-C_0.
\]
There is no interior crossing at a seam or an endpoint translate.
At an endpoint, $C$ includes half of an atom already of mass
$\zeta/2$, hence the contribution $\zeta/4$.

\begin{lemma}[Arc turn and seam count]\label{j:marked-phase}
Let $T(X)$ be the total turn of the lifted arc, and let $K$ be its
integer signed number of eastward crossings of the upper seam rays
$\{\ell M\}\times(0,\infty)$, summed over $\ell$. Then
\begin{equation}\label{j:marked-phase-identities}
 T(X)=-2\pi\bigl(C_{\rm per}(i)+C_{\rm per}(j)+2C_0\bigr),
 \qquad K=C_0+\frac{k}{2}.
\end{equation}
Consequently, for $q=e^{i\pi/4}$,
\[
 q^{4C_0}i^k(-1)^K=e^{2\pi iK}=1.
\]
\end{lemma}

\begin{proof}
For a seam or endpoint abscissa $a$, use the branch
$\arg(a-X)\in(0,2\pi)$. Each upward cut crossing
to the left of $a$ requires a correction $-2\pi$. The initial and
terminal branch arguments are respectively
\[
 \pi\bigl(1+\zeta_i\mathbf1_{i<a}\bigr),\qquad
 \pi\bigl(1-\zeta_j\mathbf1_{j+kM<a}\bigr),
\]
where an indicator is $1/2$ at equality. Thus the continuous
variation of $\arg(a-X(t))$, using one-sided endpoint limits, is
$-2\pi C(a)$.

Lift $\arg(X(t)-X(s))$ continuously on the triangle $s<t$.
On the route $(0,0)\to(0,1)\to(1,1)$, its changes are
$-2\pi C(i)$ and $-2\pi C(j+kM)$. On the diagonal its limiting
angle is the tangent angle, whose change is $T(X)$. The two routes
have the same endpoints, so
$T(X)=-2\pi(C(i)+C(j+kM))$. For $\ell\ne0$ instead lift
$\arg(X(t)+\ell M-X(s))$ on the whole square. Disjointness from
the translate makes this well defined. The diagonal angle is
constant; the same two-edge route therefore gives
\[
 C(i-\ell M)+C(j+kM+\ell M)=0.
\]
Summing these identities with the $\ell=0$ turn identity gives the
first formula in \eqref{j:marked-phase-identities}.

For a seam at $a$, let $I_a$ be the interior signed upward cut
flux to its left and let $K_a$ be the signed eastward crossings
of its upper ray. Flux from the upper-left region gives
\[
 K_a=I_a+\mathbf1_{\{\zeta_i=1,\ i<a\}}
          -\mathbf1_{\{\zeta_j=-1,\ j+kM<a\}}.
\]
Substituting the endpoint half-atoms into $C(a)$ makes this
\[
 K_a=
 C(a)+\frac{\mathbf1_{i<a}-\mathbf1_{j+kM<a}}2.
\]
An interior upward cut crossing enters that region and a
downward one leaves it; an upward start contributes $+1$, and
a downward finish contributes $-1$. Summing over $a=\ell M$
gives $K=C_0+k/2$, because
$0<i<j<M$ makes the sum of the indicator differences equal to
$k$. Here $K$ is an integer by its definition as a signed crossing
count. The displayed phase cancellation now follows by substituting
$C_0+k/2=K$.
\end{proof}

\subsection{Identification of the physical matrix element}

\begin{proof}[Proof of Proposition~\ref{j:marked-physical}]
For the physical site list, Lemma~\ref{j:shifted-nonvanishing}
gives $N\ne0$: after a common rescaling its logarithmic heights lie in
$\{0,1\}$, so their spread is at most one. Lemma~\ref{j:selected-vacuum-vectors} therefore
continues the selected vectors $R_+$ and $L_-$ componentwise to
this list, including its simultaneous coincidences. Their vacuum
components remain one, and the regular physical row matrices still
fix them.

We identify these algebraic vectors with the physical vacuum
limits. Stack the uniform-twist physical transfers, with twist
$i$ or $-i$, on the fixed cylinder modulo $P$. Untouched closed
loops cancel in the oriented sum. With any fixed rim states, the
remaining finitely many occupied rim ports are paired by oriented
arcs. The absolute weights of those arcs are bounded by products
of the period-excluding SAW sums of
Lemma~\ref{j:period-excluded}. Arcs reaching depth tending to
infinity from a rim have vanishing total weight, by the tail of
that finite sum. After discarding arcs that reach the middle row,
the bottom and top configurations occupy disjoint slabs and hence
factor. Thus, as the distance between the rims tends to infinity,
the transfer matrix factors into independent bottom and top vacuum
vectors. Their vacuum components are one. This
argument is for each fixed $P$ and needs no uniformity in its
length. Multiplying the continued fixed vectors by these large
powers identifies them, by their vacuum normalization, with the
corresponding physical vectors.

It remains to compute the finite-stack matrix element with the two
flips. On the reference geometry, the flips create two directed
arcs from $i$ to $j$, with opposite vertical signs at each
endpoint. Every other loop is disjoint from both arcs and cannot
separate the endpoints. Its interval insertion phase is therefore
one, so the usual orientation cancellation, including its seam
twist, removes it. The two arcs join to a single unoriented polygon.

For either arc, Lemma~\ref{j:marked-phase} makes its turning phase
\[
 e^{-iT(X)/4}
 =q^{2(C_{\rm per}(i)+C_{\rm per}(j)+2C_0)}.
\]
The total signed seam crossing count is $k$, and the signed upper
count is $K$. The lower twist $i$ and upper twist $-i$ thus
contribute $i^k(-1)^K$. Between the two arcs the endpoint atoms
cancel. Since the seam immediately precedes $i$, their summed
$C_{\rm per}(i)$ is zero. Physical spin counts downward crossings
positively, so the interval insertion equals
\[
 i^{\sum_{a\in A'}s_a^{\rm phys}}
 =q^{-2\sum_{\rm arcs}(C_{\rm per}(i)+C_{\rm per}(j))}.
\]
The remaining phase for each arc is
$q^{4C_0}i^k(-1)^K=1$. This holds for arbitrary integer $k$.
Each polygon has exactly the specified source-to-sink flow, so
its scalar weight is counted once and is $x^{|\lambda|}$.

The absolute arc sums above allow both rim heights to tend to
infinity in this finite-stack identity. They also prove finiteness
of the positive polygon sum. Hence
$H^{\rm flip}=H^{\rm poly}$ at the physical list.
\end{proof}

At a physical list we may henceforth write
$H_{\rm lg}=H^{\rm poly}=H^{\rm flip}$. At nearby generic lists
only $H^{\rm flip}$ is used. In particular any use of
\eqref{c7:eq-L} at a nongeneric list means the limit of its combined
expression. The next section estimates that algebraic quantity
before taking the physical limit.
